\chapter*{Предисловие}
\addcontentsline{toc}{chapter}{Предисловие}

Этот текст нужен не для того, чтобы еще раз кратко пересказать ноутбуки, а для
того, чтобы собрать из них единую исследовательскую историю. В серии
\texttt{exp2--exp9a} происходило не линейное ``улучшение метрики'', а медленное
смещение вопроса. Сначала задача выглядела как инженерная: построить эмбрион
памяти из реальных источников и проверить, что по нему вообще можно пустить
сигнал. Потом стало ясно, что сам факт работоспособности пайплайна почти ничего
не говорит о качестве retrieval. Затем выяснилось, что недостаточно просто
добавить семантический gate: на уровне топологии и структуры сходимости граф
ведет себя иначе, чем хотелось бы. После этого линия распалась на несколько
ветвей:

\begin{enumerate}[leftmargin=2em]
\item ветвь локальной физики сигнала: что удерживает входящий поток и чем
отличается узел-ответ от узла-хаба;
\item ветвь глобальных поправок: как отделить яркость, вызванную конкретным
запросом, от яркости, вызванной общей структурной славой узла;
\item ветвь переноса на реальный эмбрион: сохраняются ли выводы, полученные на
читабельных золотых мирах;
\item ветвь транспортной реконструкции: не в скоринге ли проблема, а в том, что
сигнал вообще идет не туда или не доходит до цели;
\item ветвь качества самой геометрии сходства: не слишком ли грязен косинус в
узком конусе NumberBatch-пространства;
\item ветвь реконструкции недостающих путей: можно ли достраивать граф
семантическим каналом, не превратив его в обычный векторный поиск.
\end{enumerate}

Во всех этих ветвях важно различать три уровня утверждений.

\begin{enumerate}[leftmargin=2em]
\item \textbf{Положительный результат.} Эксперимент убедительно показывает, что
некоторый механизм действительно выполняет свою локальную работу.
\item \textbf{Отрицательный результат.} Эксперимент убедительно показывает, что
некоторая популярная надежда не срабатывает.
\item \textbf{Диагностический результат.} Эксперимент не дает победителя, но
сильно сужает пространство следующих гипотез.
\end{enumerate}

Серия \texttt{exp2--exp9a} особенно богата именно третьим типом результата.
Это не недостаток. Напротив, в архитектурной работе диагностический прогресс
часто важнее преждевременной победы: он позволяет отличить настоящий предел
механизма от случайного неуспеха конкретной настройки.

\begin{takeaway}
Главная задача этого свода -- не объявить задним числом, что серия ``доказала'',
а аккуратно восстановить логику движения. В ряде мест более ценно не то, что
некоторый метод оказался лучшим, а то, что удалось честно показать: текущий тип
механизма уперся в свой естественный потолок, и дальнейшее подкручивание тех же
ручек уже не обещает качественного скачка.
\end{takeaway}

\section*{Что именно считается экспериментом в этой книге}

Здесь под ``экспериментом'' понимается не только исполненный ноутбук. В конце
серии появляются и такие ветви, которые на момент свода существуют в разном
статусе:

\begin{itemize}[leftmargin=2em]
\item полностью исполненный ноутбук с сохраненными output-ячейками и отчетом;
\item исполненный ноутбук, но без окончательного markdown-репорта;
\item проектная развилка, подробно спроектированная как следующий опыт, но еще
не прогнанная.
\end{itemize}

Это особенно важно для \texttt{exp8b}. Его нельзя молча подать как уже
завершенный опыт. Но и выбрасывать его из исследовательской линии было бы
ошибкой: он играет роль концептуального фильтра, который должен был отделить
дефекты транспортной архитектуры от дефектов самой embedding-геометрии.

\section*{Как читать этот текст}

Есть два разумных маршрута чтения.

\begin{enumerate}[leftmargin=2em]
\item \textbf{Линейный маршрут.} Читать подряд, как историю исследования. Тогда
лучше начинать с введения и не перепрыгивать через главы, потому что поздние
серии (\texttt{7x--9a}) содержательно опираются на различия, выработанные в
\texttt{6x}.
\item \textbf{Маршрут по вопросам.} Если интересует одна конкретная тема --
например, PageRank-поправка, capacity, overlap cue-fields или semantic
proximity edges, -- можно идти сразу в соответствующую серию. Но в этом случае
нужно помнить, что поздние главы используют терминологию, введенную в первых
двух главах.
\end{enumerate}

\section*{Стилистическое замечание}

Язык текста намеренно педагогичен, но не упрощен до уровня популяризации.
Читатель здесь предполагается математически и алгоритмически подготовленным.
Поэтому базовые понятия вроде графа, случайного блуждания, косинусной
схожести или стационарного распределения не разжевываются с нуля; вместо этого
основной акцент сделан на том, как именно эти объекты начинают работать внутри
архитектуры памяти и почему в одном месте они помогают, а в другом создают
иллюзию прогресса.
